\begin{abstract}
The intricacies and formal structures inherent in mathematical reasoning make it particularly challenging for language models to master. In this study, we present DeepSeekMath~7B, which is developed by further pre-training DeepSeek-Coder-Base-v1.5 7B on 120B math-centric tokens drawn from Common Crawl, as well as supplementary natural language and programming data. DeepSeekMath~7B attains an outstanding 51.7\% accuracy on the MATH benchmark at the competition level, without employing external solvers or ensemble methods, closely rivaling the scores achieved by Gemini-Ultra and GPT-4. With self-consistency over 64 samples, DeepSeekMath~7B increases its MATH benchmark score to 60.9\%. The advancements in mathematical reasoning achieved by DeepSeekMath~are underpinned by two principal innovations: firstly, a carefully constructed pipeline for selecting and curating public web data enables the exploitation of large-scale open resources; secondly, the introduction of Group Relative Policy Optimization (GRPO)—an adaptation of Proximal Policy Optimization (PPO)—boosts mathematical problem-solving proficiency while reducing PPO’s memory requirements.
\end{abstract}

\section{Introduction}

The advent of large language models (LLMs) has dramatically transformed the field of mathematical reasoning within artificial intelligence, enabling notable progress across quantitative reasoning evaluations \citep{MATH} and tests of geometric problem-solving \citep{trinh2024solving}. These systems have also become valuable tools for aiding human problem-solvers in navigating intricate mathematical tasks \citep{tao}. Despite these advancements, premier models like GPT-4 \citep{gpt4} and Gemini-Ultra \citep{gemini} remain inaccessible to the public, leaving existing open-source alternatives substantially lagging in performance.

This work presents DeepSeekMath, a specialized language model that markedly surpasses open-source alternatives in mathematical reasoning, and attains performance close to that of GPT-4 across standard academic assessments. Central to this achievement is the construction of the DeepSeekMath Corpus, a comprehensive, high-quality pre-training resource containing 120B math-specific tokens. This dataset is derived from Common Crawl (CC) via a fastText-based classification system \citep{joulin2016fasttext}. In the preliminary phase, the classifier is trained with positive samples from OpenWebMath \citep{openwebmath}, supplemented by a varied array of unrelated web content as negative samples. The trained classifier is subsequently utilized to identify further positive examples from CC, which are then subjected to human annotation for increased reliability. The classifier undergoes retraining with the expanded and refined dataset to enhance its efficacy. Empirical evaluation demonstrates that the resultant corpus is of substantial quality, with DeepSeekMath-Base 7B achieving 64.2\% on GSM8K \citep{gsm8k} and 36.2\% on the competitive-level MATH dataset \citep{MATH}, exceeding the performance of Minerva 540B \citep{minerva}. Furthermore, the multilingual aspect of the DeepSeekMath Corpus yields gains on Chinese mathematical benchmarks \citep{wei2023cmath,agieval}. We posit that our methodologies in constructing mathematical datasets represent a foundational step for subsequent research, acknowledging considerable scope for future enhancement.

DeepSeekMath-Base is initialized from DeepSeek-Coder-Base-v1.5 7B \citep{deepseek-coder}, as empirical evidence suggests that commencing training from a code-specialized foundation yields superior outcomes relative to a standard large language model (LLM). Notably, we observe that training with mathematical datasets enhances not only the model's mathematical proficiency but also its performance on broader benchmarks, including MMLU \citep{mmlu} and BBH \citep{bbh}, demonstrating that such training bolsters general reasoning skills alongside domain-specific capability.

Following the pre-training phase, we subject DeepSeekMath-Base to instruction tuning on mathematical reasoning data, which encompasses chain-of-thought approaches \citep{cot}, program-of-thought exemplars \citep{pot,pal}, and tool-assisted reasoning \citep{tora}. The model produced by this tuning, DeepSeekMath-Instruct 7B, surpasses all other 7B models and achieves results comparable to open-source instruction-tuned models of 70B parameters.

Additionally, we propose Group Relative Policy Optimization (GRPO), a modified reinforcement learning (RL) approach based on Proximal Policy Optimization (PPO) \citep{schulman2017proximal}. Unlike PPO, GRPO dispenses with the critic network and instead derives baseline estimates from group-based scores, thus notably cutting down on computational demands during training. Using only a selected subset of English instruction-tuning data, GRPO delivers marked improvements over the robust DeepSeekMath-Instruct baseline. This is evidenced by gains in both familiar domains (GSM8K: 82.9\% $\rightarrow$ 88.2\%, MATH: 46.8\% $\rightarrow$ 51.7\%) and out-of-distribution math benchmarks (e.g., CMATH: 84.6\% $\rightarrow$ 88.8\%) during the RL stage. Moreover, we present a unified framework to analyze methodologies such as Rejection Sampling Fine-Tuning (RFT) \citep{yuan2023scaling}, Direct Preference Optimization (DPO) \citep{dpo}, PPO, and GRPO. This common perspective reveals that these techniques can be viewed as manifestations of direct or simplified RL strategies. We conduct thorough empirical evaluations—including comparisons between online and offline training regimes, outcome versus process-based supervision, and single-step versus iterative RL—to identify and assess critical facets of this paradigm. Finally, we elucidate the mechanisms through which our RL strategies enhance instruction-tuned models and outline future avenues for developing even more effective RL methods grounded in this unified framework.

\subsection{Contributions}
This work presents advances in large-scale math-focused pre-training and investigates reinforcement learning methods through empirical analysis.

\textbf{Math Pre-Training at Scale}

\begin{itemize}[topsep=0pt]
    \item We demonstrate that the open-source Common Crawl repository is a rich resource for mathematical data. Utilizing a carefully engineered data filtering pipeline, we assemble the DeepSeekMath Corpus—comprising 120B tokens sourced from web pages with substantial mathematical content. This dataset substantially exceeds the scale of those previously assembled for Minerva \citep{minerva} and OpenWebMath \citep{openwebmath}, being nearly seven times and nine times larger, respectively.

    \item Our base model, DeepSeekMath-Base 7B, attains performance on par with the much larger Minerva 540B model \citep{minerva}, suggesting that parameter count alone does not dictate proficiency in mathematical reasoning. Instead, pre-training on a carefully curated, high-quality mathematical dataset enables more compact models to achieve competitive results.

    \item The study provides insights into the role of code pre-training as a precursor to math pre-training. We find that code pre-training enhances a model's ability to solve mathematical questions, regardless of whether external tools are used, lending support to the hypothesis that code pre-training bolsters reasoning skills, especially in mathematical domains.

    \item Contrary to common practice, incorporating arXiv paper data, although widespread in mathematics-related works, does not result in observable improvements across the various mathematical benchmarks evaluated in our experiments.
\end{itemize}

\textbf{Exploration and Analysis of Reinforcement Learning}

\begin{itemize}[topsep=0pt]
    \item We present Group Relative Policy Optimization (GRPO), a novel reinforcement learning approach that is both resource-efficient and robust. Rather than utilizing a critic network, GRPO estimates the baseline using group-level scoring, which markedly lowers the computational requirements relative to Proximal Policy Optimization (PPO).
    \item Our experiments show that GRPO yields considerable improvements in the instruction-tuned DeepSeekMath-Instruct model by relying exclusively on instruction-tuning datasets. We also report that, through reinforcement learning, the model displays improved generalization to out-of-domain tasks.
    \item We offer a comprehensive framework for interpreting a range of methods, including RFT, DPO, PPO, and GRPO. Additionally, we carry out in-depth analyses through a variety of experiments—such as contrasting online and offline training regimes, evaluating outcome-based versus process-oriented supervision, and comparing single-turn with iterative reinforcement learning—to systematically examine the critical components underpinning this framework.
    \item Building upon our integrative framework, we investigate factors contributing to reinforcement learning's efficacy for language models, and outline several promising avenues for advancing reinforcement learning approaches with LLMs.
\end{itemize}

\subsection{Summary of Evaluations and Metrics}

\begin{itemize}[topsep=0pt]
    \item \textbf{English and Chinese Mathematical Reasoning}: Our evaluation encompasses an extensive suite of English and Chinese benchmarks, featuring mathematical questions ranging from elementary to university levels. English benchmarks surveyed include GSM8K \citep{gsm8k}, MATH \citep{MATH}, SAT \citep{llemma}, OCW Courses \citep{minerva}, and MMLU-STEM \citep{mmlu}, while the Chinese set consists of MGSM-zh \citep{mgsm}, CMATH \citep{wei2023cmath}, Gaokao-MathCloze \citep{agieval}, and Gaokao-MathQA \citep{agieval}. Assessments focus on models' capability to provide coherent written solutions independently and their competence in solving tasks programmatically using Python.

    In English language tasks, DeepSeekMath-Base shows performance on par with the proprietary Minerva 540B \citep{minerva}, and outperforms all released base models such as Mistral 7B \citep{mistral} and Llemma-34B \citep{llemma}, with or without targeted math pre-training, frequently by notable margins. Particularly in Chinese mathematics benchmarks, DeepSeekMath-Base excels, attributed to our inclusion of high-quality, diverse language pre-training data, in contrast with prior efforts \citep{minerva,llemma} that restrict themselves to English sources. Upon further refinement with instruction tuning and reinforcement learning, both DeepSeekMath-Instruct and DeepSeekMath-RL set new open-source standards, with DeepSeekMath-RL being the first to reach above 50\% accuracy on the advanced MATH competition dataset.
\end{itemize}

\item \textbf{Formal Mathematics}: We assess DeepSeekMath-Base using the informal-to-formal theorem proving task introduced in \citep{dsp_proof} on the miniF2F benchmark \citep{minif2f}, employing Isabelle \citep{isabelle} as the designated proof assistant. Our results show that DeepSeekMath-Base achieves robust few-shot performance in the autoformalization task.

\item \textbf{Natural Language Understanding, Reasoning, and Code}: To form a holistic view of the model's abilities in general language understanding, reasoning, and programming, DeepSeekMath-Base is benchmarked on Massive Multitask Language Understanding (MMLU) \citep{mmlu}, a suite of 57 diverse multiple-choice problems; BIG-Bench Hard (BBH) \citep{bbh}, a collection of 23 complex tasks with an emphasis on multi-step reasoning; and the code-oriented HumanEval \citep{codex} and MBPP \citep{mbpp} datasets. We observe that math-specific pre-training yields improvements not only in language comprehension but also in reasoning tasks.

\section{Math Pre-Training}

\subsection{Data Collection and Decontamination}
This section details our methodology for constructing the DeepSeekMath Corpus utilizing the Common Crawl dataset. Illustrated in Figure \ref{fig:data_collect}, the pipeline follows an iterative strategy to assemble a vast corpus of mathematical text, beginning with a small, curated seed dataset. While our focus is mathematics, this approach is extendable to other fields, such as programming.

Initially, we select OpenWebMath \citep{openwebmath}, which offers a reliable foundation of high-quality math web texts, as the seed dataset. Leveraging this corpus, a fastText model \citep{joulin2016fasttext} is trained to identify web pages resembling those in OpenWebMath. For training, we sample 500,000 positive examples from the seed set and an equal number of negatives from Common Crawl. Model training is carried out using an open-source toolkit\footnote{\url{https://fasttext.cc}}, with the following hyperparameters: vector size of 256, learning rate of 0.1, word n-gram size up to 3, minimum word frequency of 3, and 3 training epochs. To pare down the size of Common Crawl, we apply URL-based and near-deduplication methods, thus reducing it to 40B unique HTML web pages. Next, candidate mathematical pages are retrieved from this deduplicated set using the fastText classifier. To maintain quality, retrieved pages are ranked by their fastText scores, retaining only the highest-ranking entries. The quantity of retained data is determined through pre-training on subsets of the top 40B, 80B, 120B, and 160B tokens; in our initial round, the top 40B tokens are selected.

Following the initial round of data gathering, a significant portion of mathematical web pages remains uncaptured, primarily due to insufficient diversity within the set of positive examples used to train the fastText model. To address this, we seek out further mathematical web sources to augment the initial seed dataset, thereby enhancing the fastText model's effectiveness. Our process starts by segmenting the full Common Crawl dataset into separate domains, where each domain is composed of web pages sharing a common base URL. For every domain, we determine the proportion of web pages that were collected during the first round. Those domains from which more than 10\% of the web pages were gathered are tagged as math-oriented domains (for example, \url{mathoverflow.net}).

Next, within these detected domains, we manually review and annotate URLs known to host mathematical material (such as \url{mathoverflow.net/questions}). Uncollected web pages associated with these URLs are incorporated into the expanding seed corpus. Through this method, we compile a broader range of positive examples and retrain the fastText model, which, in turn, enhances its capacity to identify mathematical web data in future iterations. This iterative approach is repeated four times, resulting in a final set of 35.5 million mathematical web pages, amounting to 120 billion tokens. Notably, by the fourth round, we observe that approximately 98\% of the data identified has already been retrieved in the previous iteration, prompting us to halt further data acquisition.

To mitigate the risk of benchmark data leakage, we adopt the filtering strategy presented in \cite{deepseek-coder} and eliminate web pages containing questions or answers from widely-used English mathematical benchmarks (GSM8K \citep{gsm8k}, MATH \citep{MATH}) as well as Chinese benchmarks (CMATH \citep{wei2023cmath}, AGIEval \citep{agieval}). Our approach entails removing from the math training corpus any text fragment that contains a 10-gram exact match with any sequence from these benchmarks. For benchmark texts shorter than 10 grams but consisting of at least 3 grams, we rely on precise matching to exclude contaminated web pages.

\subsection{Validating the Quality of the DeepSeekMath~Corpus}

To assess the efficacy of the DeepSeekMath~Corpus relative to other recent large-scale mathematics corpora, we conduct pre-training experiments comparing it to:

\begin{itemize}[topsep=0pt]
    \item \textbf{MathPile} \citep{mathpile}: A diverse corpus totaling 8.9B tokens, drawn from sources including textbooks, Wikipedia, ProofWiki, CommonCrawl, StackExchange, and arXiv, with over 85\% of its data originating from arXiv;
    \item \textbf{OpenWebMath} \citep{openwebmath}: Constructed from CommonCrawl and filtered for mathematics-related content, comprising 13.6B tokens;
    \item \textbf{Proof-Pile-2} \citep{llemma}: Composed of OpenWebMath, AlgebraicStack (10.3B tokens of mathematical code), and arXiv manuscripts (28.0B tokens). When utilizing Proof-Pile-2 for experiments, we adhere to the arXiv:Web:Code ratio of 2:4:1 as outlined in \cite{llemma}.
\end{itemize}

\subsubsection{Training Setting}
\label{sec:quality-policy}
We perform mathematical domain adaptation on a general-purpose pre-trained language model consisting of 1.3B parameters, which is architecturally aligned with the DeepSeek LLMs \citep{deepseek-llm}, referred to herein as DeepSeek-LLM 1.3B. Each mathematical dataset is used independently to train a separate model for a total of 150B tokens. All training runs are executed with the efficient and lightweight HAI-LLM \citep{haillm} infrastructure. Adhering to the methodology used by DeepSeek LLMs, we employ the AdamW optimizer \citep{adamW}, setting $\beta_1=0.9$, $\beta_2=0.95$, and $\mathrm{weight\_decay}=0.1$. The learning rate schedule follows a multi-phase strategy: the rate peaks following a 2,000-step warmup, then drops to 31.6\% at the 80\% mark of training completion, and further decays to 10.0\% by the 90\% point. The upper bound of the learning rate is established at 5.3e-4, and model optimization is performed with a batch size of 4 million tokens and a 4,000-token context window.

\subsubsection{Evaluation Results}

\textbf{The DeepSeekMath~Corpus exhibits outstanding quality, includes diverse multilingual mathematical materials, and surpasses all others in scale.}

\begin{itemize}[topsep=0pt]
    \item \textbf{High-quality}:
    We assess model capability on eight mathematical evaluation benchmarks utilizing few-shot chain-of-thought prompting \cite{cot}. According to the results in Table \ref{tab:corpora_comparison}, the model trained with the DeepSeekMath~Corpus achieves a clear performance advantage. Additionally, Figure \ref{fig:corpora_comparisons} highlights that, at the 50B token mark (corresponding to one full epoch of Proof-Pile-2), the DeepSeekMath-trained model notably outperforms the Proof-Pile-2-trained counterpart, indicating a superior average quality of the DeepSeekMath Corpus.
    \item \textbf{Multilingual}:
    DeepSeekMath~Corpus contains resources in a variety of languages, with English and Chinese comprising the majority. Table \ref{tab:corpora_comparison} illustrates that fine-tuning on the DeepSeekMath~Corpus leads to improved mathematical reasoning in both English and Chinese. By contrast, previous mathematical datasets, which are mostly English-dominant, show marginal benefits and sometimes negative transfer to Chinese mathematical tasks.
    \item \textbf{Large-scale}:
    DeepSeekMath~Corpus is substantially larger than any previously available mathematical corpus. In Figure \ref{fig:corpora_comparisons}, training DeepSeek-LLM 1.3B on the DeepSeekMath~Corpus produces a steeper and more sustained learning progression. On the other hand, alternative baseline corpora are much smaller and are seen multiple times during training, resulting in earlier saturation of performance gains.
\end{itemize}

\subsection{Training and Evaluating DeepSeekMath-Base 7B}

Here, we present DeepSeekMath-Base 7B, a foundational model designed to demonstrate strong mathematical reasoning skills. This model is initialized using DeepSeek-Coder-Base-v1.5 7B \citep{deepseek-coder} and is further trained on 500B tokens. The composition of the training corpus is as follows: 56\% DeepSeekMath Corpus, 4\% AlgebraicStack, 10\% arXiv, 20\% Github code, and the last 10\% consists of natural language data from Common Crawl in both English and Chinese. The majority of the training parameters align with those detailed in Section \ref{sec:quality-policy}, with exceptions including setting the peak learning rate to 4.2e-4 and employing a batch size of 10M tokens.

To thoroughly evaluate DeepSeekMath-Base 7B, we analyze its proficiency in generating self-sufficient mathematical solutions independent of auxiliary tools, its capability in tool-augmented problem solving, and its competency in formal theorem verification. In addition to mathematical problem solving, we also report a broader assessment of the base model’s aptitude in natural language understanding, logical reasoning, and programming abilities.

\paragraph{Mathematical Problem Solving with Step-by-Step Reasoning}

The effectiveness of DeepSeekMath-Base in resolving mathematical tasks is measured using few-shot chain-of-thought prompting \citep{cot}, spanning eight distinct benchmarks in both English and Chinese. The benchmarks cover a range of quantitative reasoning tasks—such as GSM8K \citep{gsm8k}, MATH \citep{MATH}, and CMATH \citep{wei2023cmath}—as well as multiple-choice assessments, including MMLU-STEM \citep{mmlu} and Gaokao-MathQA \citep{agieval}, thus representing a wide array of mathematical domains from elementary to university-level challenges.

As indicated in Table \ref{tab:base_math_cot}, among open-source base models—including the commonly utilized Mistral 7B \citep{mistral} and the newly introduced Llemma 34B \citep{llemma} trained on Proof-Pile-2 \citep{llemma}—DeepSeekMath-Base 7B achieves superior results on all eight evaluated benchmarks. Remarkably, for the challenging MATH dataset, DeepSeekMath-Base delivers an absolute gain exceeding 10\% over other open-source base models, and even outperforms Minerva 540B \citep{minerva}, a proprietary base model with a parameter count 77 times larger, which is itself built upon PaLM \citep{palm} and additionally refined with mathematical content.

\paragraph{Mathematical Problem Solving with Tool Use} To measure the capacity for program-assisted mathematical reasoning, we assess few-shot program-of-thought prompting \citep{pot,pal} on GSM8K and MATH. Here, models are given prompts that require solving problems through constructing Python scripts—permitting the use of libraries such as \textit{math} and \textit{sympy} for advanced calculations—and the resulting output from program execution is regarded as the model's answer. According to Table \ref{tab:base_math_pot_proof}, DeepSeekMath-Base 7B achieves results surpassing those of the previous state-of-the-art, Llemma 34B.

\paragraph{Formal Mathematics}

Automating formal proof generation offers improvements in correctness, reliability, and efficiency, which has attracted growing research attention. DeepSeekMath-Base 7B is tested on the informal-to-formal proof task from \citep{dsp_proof}, which involves producing a formal proof from an informal premise, a formal translation of that statement, and an informal proof outline. Using the miniF2F benchmark \citep{minif2f}, which addresses Olympiad-level formal mathematics, formal proofs in Isabelle are generated with few-shot prompting for each example. As described in \cite{dsp_proof}, we have models create proof outlines and use the standard automated theorem prover Sledgehammer \citep{sledgehammer} to complete any remaining proof details. Results in Table \ref{tab:base_math_pot_proof} demonstrate that DeepSeekMath-Base 7B achieves robust performance in the autoformalization of proofs.

\paragraph{Natural Language Understanding, Reasoning, and Code}

We assess the model’s ability in natural language understanding via MMLU \citep{mmlu}, its reasoning skills with BBH \citep{bbh}, and programming proficiency on HumanEval \citep{codex} and MBPP \citep{mbpp}. Table \ref{tab:understanding_reasoning_code} shows that DeepSeekMath-Base 7B registers notable improvements on MMLU and BBH compared to its predecessor DeepSeek-Coder-Base-v1.5 \citep{deepseek-coder}, underscoring the value of mathematical pretraining for language comprehension and reasoning. Furthermore, by continuing pretraining with code-related tokens, DeepSeekMath-Base 7B preserves the coding abilities of DeepSeek-Coder-Base-v1.5 on both programming benchmarks. Overall, DeepSeekMath-Base 7B demonstrates a clear performance advantage over the general-purpose Mistral 7B model \citep{mistral} across all three tasks involving reasoning and code.

\section{Supervised Fine-Tuning}

\subsection{SFT Data Curation}

To prepare a comprehensive dataset for mathematical instruction tuning, we gathered problems in both English and Chinese, spanning a wide array of mathematical disciplines and difficulty levels. Each problem is matched with corresponding solutions employing chain-of-thought (CoT) \citep{cot}, program-of-thought (PoT) \citep{pot,pal}, or tool-integrated reasoning approaches \citep{tora}. In total, the dataset contains 776K training instances.

\begin{itemize}[topsep=0pt]
    \item \textbf{English mathematical datasets}: We augment GSM8K and MATH with tool-assisted annotated answers, supplement these with select items from MathInstruct \citep{MathInstruct}, and utilize the Lila-OOD training split \citep{lila} where solutions are provided via CoT or PoT. Our English data covers mathematics fields including, but not limited to, algebra, probability, number theory, calculus, and geometry.
    \item \textbf{Chinese mathematical datasets}: We curate Chinese K-12 mathematics questions covering 76 distinct sub-topics, such as linear equations, where each is furnished with both chain-of-thought and tool-integrated explanations.
\end{itemize}

\subsection{Training and Evaluating DeepSeekMath-Instruct 7B}

This section details the development of DeepSeekMath-Instruct 7B, which is refined with mathematical instruction tuning on top of DeepSeekMath-Base. During training, data samples are concatenated at random, up to a context length ceiling of 4K tokens. The model is trained for 500 steps, employing a batch size of 256 and a fixed learning rate of 5e-5.

For performance assessment, we examine mathematical capabilities in settings both with and without tool assistance, applying four quantitative reasoning evaluation sets in both English and Chinese. Comparisons are made with other state-of-the-art models:

\begin{itemize}[topsep=0pt]
    \item \textbf{Closed-source models} considered for comparison include: (1) models from the GPT series, such as GPT-4 \citep{gpt4} and GPT-4 Code Interpreter~\footnote{\url{https://openai.com/blog/chatgpt-plugins\#\#code-interpreter}}, (2) Gemini Ultra and Pro \citep{gemini}, (3) Inflection-2 \citep{inflection-2}, (4) Grok-1~\footnote{\url{https://x.ai/model-card}}, in addition to several newly available Chinese models, including (5) Baichuan-3~\footnote{\url{https://www.baichuan-ai.com}}, and (6) the newest member of the GLM line, GLM-4~\footnote{\url{https://open.bigmodel.cn/dev/api\#glm-4}} from the GLM family \citep{glm}.
\end{itemize}

These models are designed for broad applications and have typically been subjected to various alignment processes.
\item \textbf{Open-source models} consist of: (1) DeepSeek-LLM-Chat 67B \citep{deepseek-llm}, (2) Qwen 72B \citep{qwen}, (3) SeaLLM-v2 7B \citep{seallm}, and (4) ChatGLM3 6B \citep{chatglm3} as general-purpose systems, alongside math-enhanced architectures such as (5) InternLM2-Math 20B~\footnote{\url{https://github.com/InternLM/InternLM-Math}}—an extension of InternLM2 with dedicated mathematical training and subsequent instruction tuning, (6) Math-Shepherd-Mistral 7B, which utilizes PPO optimization \citep{schulman2017proximal} with Mistral 7B \citep{mistral} through a reward model supervised by intermediate solutions, (7) the WizardMath series \citep{wizardmath}, which boosts the mathematical reasoning ability of Mistral 7B and Llama-2 70B \citep{llama2} by applying evolve-instruct (leveraging AI-generated instruction variations) and PPO, training on tasks primarily from GSM8K and MATH, (8) MetaMath 70B \citep{metamath}—Llama-2 70B further fine-tuned with an expanded GSM8K and MATH dataset, (9) ToRA 34B \cite{tora}—a version of CodeLlama 34B specialized for mathematical problem solving with tool integration, and (10) MAmmoTH 70B \citep{MathInstruct}, where Llama-2 70B undergoes MathInstruct-based instruction tuning.
\end{itemize}

As depicted in Table \ref{tab:sft_rl_math}, when tool use is not permitted during evaluation, DeepSeekMath-Instruct 7B delivers robust performance on stepwise reasoning tasks. On the competitive MATH dataset in particular, our model exceeds all other open-source alternatives and outperforms most proprietary solutions (such as Inflection-2 and Gemini Pro) by margins of at least 9 percentage points. This advantage persists even over models with considerably more parameters (e.g., Qwen 72B) or those trained specifically with math-focused reinforcement learning strategies (such as WizardMath-v1.1 7B). Although DeepSeekMath-Instruct performs on par with leading Chinese proprietary models like GLM-4 and Baichuan-3 on MATH, it does not match the performance levels of GPT-4 or Gemini Ultra.

When the evaluation protocol permits combining natural language and programmatic tool use for problem-solving, DeepSeekMath-Instruct 7B attains an accuracy close to 60\% on the MATH dataset—outperforming every open-source peer. For other benchmarks, our model’s results are comparable to those of DeepSeek-LLM-Chat 67B, the previous best model in the class, despite being only one-tenth the size.

\section{Reinforcement Learning}

\subsection{Group Relative Policy Optimization} Reinforcement learning (RL) has proven particularly effective for advancing the mathematical reasoning skills of LLMs after completion of the Supervised Fine-Tuning (SFT) phase \citep{wang2023math,wizardmath}. This section describes our proposed RL approach, Group Relative Policy Optimization (GRPO), which is both efficient and impactful.

\subsubsection{From PPO to GRPO} Proximal Policy Optimization (PPO) \citep{schulman2017proximal}, an actor-critic RL approach, is a staple for RL fine-tuning of LLMs \citep{ouyang2022training}. It refines models by maximizing the following surrogate objective:

\begin{equation}
\footnotesize
    \mathcal{J}_{PPO}(\theta) = \mathbb{E}{[q \sim P(Q), o \sim \pi_{\theta_{old}}(O|q)]} \frac{1}{|o|} \sum_{t=1}^{|o|} \min \left[ \frac{\pi_\theta(o_{t} | q, o_{<t})}{\pi_{\theta_{old}}(o_{t} | q, o_{<t})} A_{t}, \text{clip} \left( \frac{\pi_\theta(o_{t} | q, o_{<t})}{\pi_{\theta_{old}}(o_{t} | q, o_{<t})}, 1 - \varepsilon,

Here, $\pi_{\theta}$ and $\pi_{\theta_{old}}$ represent the current and prior policy networks, respectively, while $q$ and $o$ denote questions and sampled outputs drawn from the question dataset and the preceding policy $\pi_{\theta_{old}}$. The hyper-parameter $\varepsilon$ is incorporated as part of the clipping mechanism in PPO to promote training stability. The advantage term, $A_t$, is estimated using Generalized Advantage Estimation (GAE) \citep{gae}, which leverages reward signals $\{r_{\ge t}\}$ along with a learned value function $V_{\psi}$. Consequently, PPO requires simultaneous optimization of both a value function and the policy network. To counteract reward model overfitting, it is customary to integrate a per-token Kullback-Leibler (KL) penalty—relative to a reference model—into the reward at each generation step, as proposed in \citep{ouyang2022training}, leading to the formulation:

\begin{equation}
    r_{t} = r_\varphi(q, o_{\le t}) - \beta \log\frac{\pi_{\theta}(o_{t}|q, o_{<t})}{\pi_{ref}(o_{t}|q, o_{<t})},
\label{eq:PPO-reward}
\end{equation}

where $r_\varphi$ is the reward model, $\pi_{ref}$ indicates the reference policy (typically the initial SFT model), and $\beta$ is the coefficient scaling the KL penalty.

Training a value function in PPO often requires maintaining an additional model with similar computational and memory demands as the policy network. During reinforcement learning, the value function serves as a variance-reducing baseline for advantage computation. In the context of LLMs, where rewards are generally assigned to only the final token by the reward model, learning an accurate value function at every token can be especially challenging. To alleviate these concerns, we introduce Group Relative Policy Optimization (GRPO), as depicted in Figure \ref{fig:grpo}. GRPO eliminates the necessity for a separately trained value function by utilizing the mean reward from a set of outputs sampled for the same prompt as a baseline. Specifically, for a question $q$, we generate a batch of responses $\{o_1, o_2, \cdots, o_G\}$ using $\pi_{\theta_{old}}$, then update the policy by maximizing the following expectation:

\begin{equation}
\footnotesize
\begin{split}
    \mathcal{J}_{GRPO}(\theta) &= \mathbb{E}{[q \sim P(Q), \{o_i\}_{i=1}^G \sim \pi_{\theta_{old}}(O|q)]}  \\
    & \frac{1}{G}\sum_{i=1}^G\frac{1}{|o_i|} \sum_{t=1}^{|o_i|} \left\{ \min \left[ \frac{\pi_\theta(o_{i,t} | q, o_{i,<t})}{\pi_{\theta_{old}}(o_{i,t} | q, o_{i,<t})} \hat{A}_{i,t}, \text{clip} \left( \frac{\pi_\theta(o_{i,t} | q, o_{i,<t})}{\pi_{\theta_{old}}(o_{i,t} | q, o_{i,<t})}, 1 - \varepsilon, 1 + \varepsilon \right)  \hat{A}_{i,t} \right] - \beta \mathbb{D}_{KL}\left[\pi_{\theta} || \pi_{ref}\right]\right\} ,
\end{split}
\label{eq:GRPO-obj}
\end{equation}

where both $\varepsilon$ and $\beta$ are hyper-parameters, and $\hat{A}_{i,t}$ is the advantage, determined by comparing only the rewards among outputs within the same group (detailed further in subsequent sections). This group-relative approach naturally complements the inherent comparison structure of reward model training data, which generally consists of output pairings for shared questions. Moreover, unlike the conventional method of including the KL penalty within the reward computation, GRPO directly integrates the KL divergence between the current policy and the reference policy into the loss term, thus keeping the calculation of $\hat{A}_{i,t}$ straightforward. Additionally, instead of the KL penalty specified in (\ref{eq:PPO-reward}), GRPO estimates the KL divergence by the following unbiased estimator \citep{kl_approx}:

\begin{equation}
\small
    \mathbb{D}_{KL}\left[\pi

\begin{algorithm}[t]
  \small
  \caption{Iterative Group Relative Policy Optimization}
  \textbf{Input} initial policy model $\pi_{\theta_{\text{init}}}$; reward models $r_\varphi$; task prompts $\mathcal{D}$; 
  hyperparameters $\varepsilon$, $\beta$, $\mu$
  \begin{algorithmic}[1]
    \State Initialize the policy model: $\pi_\theta \leftarrow \pi_{\theta_{\text{init}}}$
    \For{iteration = 1, \dots, I}
       \State Assign the current policy as the reference model: $\pi_{ref} \leftarrow \pi_{\theta}$
      \For{step = 1, \dots, M}
      \State Select a mini-batch $\mathcal{D}_b$ sampled from $\mathcal{D}$
      \State Set the previous policy model: $\pi_{\theta_{old}} \leftarrow \pi_{\theta}$ 
      \State For each query $q$ in $\mathcal{D}_b$, generate $G$ samples $\{o_i\}_{i=1}^G \sim \pi_{\theta_{old}} (\cdot \mid q)$
      \State Evaluate the sampled outputs using $r_{\varphi}$, obtaining $\{r_i\}_{i=1}^{G}$
      \State Calculate group relative advantages $\hat{A}_{i,t}$ for the $t$-th token in each output $o_i$
      \For{GRPO iteration = 1, \dots, $\mu$}
        \State Perform a policy update on $\pi_{\theta}$ by maximizing the GRPO loss (see Equation \ref{eq:GC-GRPO})
      \EndFor
    \EndFor \State Refine $r_\varphi$ with ongoing training employing a replay strategy. 
    \EndFor 
  \end{algorithmic}
  \textbf{Output} $\pi_\theta$
  \label{alg:iter-grpo}
\end{algorithm}

\subsubsection{Outcome Supervision RL with GRPO} More specifically, for each query $q$, a collection of outputs $\{o_1, o_2, \cdots, o_G\}$ is generated from the previous policy $\pi_{\theta_{old}}$. Each output is then evaluated by the reward model, resulting in a set of $G$ reward values $\mathbf{r}=\{r_1, r_2, \cdots, r_G\}$. These rewards are standardized by removing the group mean and dividing by the group standard deviation. With outcome supervision, the normalized reward corresponding to each output $o_i$ is assigned as the advantage for every token within $o_i$, that is, $\hat{A}_{i, t} = \widetilde{r}_i = \frac{r_i- {\rm mean}(\mathbf{r})}{{\rm std}(\mathbf{r})}$, and the policy parameters are optimized to maximize the objective expressed in equation (\ref{eq:GRPO-obj}).

\subsubsection{Process Supervision RL with GRPO} While outcome supervision delivers a terminal reward for each output, such feedback may be limited and inefficient, particularly in intricate mathematical tasks. Building on \cite{wang2023math}, we also investigate process supervision, which evaluates each step of reasoning by providing intermediate rewards. Formally, for a given query $q$ and sampled responses $\{o_1, o_2, \cdots, o_G\}$, a process reward model evaluates every step, resulting in the reward collections: $\mathbf{R} = \{ \{r_1^{index(1)},\cdots,r_1^{index(K_1)}\}, \cdots,  \{r_G^{index(1)},\cdots,r_G^{index(K_G)}\} \}$, where $index(j)$ identifies the end token of step $j$ and $K_i$ is the step count in $o_i$. Each step's reward is then standardized as $\widetilde{r}_i^{index(j)} = \frac{r_i^{index(j)} - {\rm mean(\mathbf{R})}}{{\rm std(\mathbf{R})}}$. Process supervision subsequently determines token-level advantages by summing the normalized future rewards: $\hat{A}_{i, t} = \sum_{index(j) \ge t} \widetilde{r}_i^{index(j)}$, and policy learning is performed by maximizing the objective in equation (\ref{eq:GRPO-obj}).

\subsubsection{Iterative RL with GRPO} During the reinforcement learning process, the original reward model may no longer provide adequate supervision for the evolving policy model. To address this, we implement iterative RL using GRPO. As described in Algorithm \ref{alg:iter-grpo}, iterative GRPO involves generating updated datasets for the reward model using outputs sampled from the latest policy model. The reward model is then retrained with a replay buffer mechanism that mixes in 10\% of prior data. Subsequently, we update the reference model to match the latest policy model and continue training the policy model, this time guided by the updated reward model.

\subsection{Training and Evaluating DeepSeekMath-RL}

Our reinforcement learning experiments utilize DeepSeekMath-Instruct 7B as the base policy model. The RL training dataset comprises chain-of-thought-style items related to GSM8K and MATH benchmarks, selected from the SFT data and totaling approximately 144K items. We intentionally omit other SFT questions to isolate the effects of RL on evaluation datasets that are not present during RL. For constructing reward model training data, we follow the methodology outlined in \citep{wang2023math}. The initial reward model is based on DeepSeekMath-Base 7B and is optimized using a learning rate of 2e-5. For the policy model during GRPO training, the learning rate is set to 1e-6, and the KL coefficient is 0.04. For each prompt, 64 candidate outputs are sampled. The maximum sequence length is capped at 1024, with a training batch size of 1024. Each exploration phase is followed by a single policy update. DeepSeekMath-RL 7B is assessed on the same set of benchmarks as DeepSeekMath-Instruct 7B. Here, GSM8K and MATH with chain-of-thought responses are considered in-domain, while all remaining benchmarks are categorized as out-of-domain tasks.

Table \ref{tab:sft_rl_math} presents a comparative analysis of both open- and closed-source models employing chain-of-thought and tool-integrated reasoning across English and Chinese datasets. Our results reveal the following insights: 1) With chain-of-thought prompting, DeepSeekMath-RL 7B achieves accuracy rates of 88.2\% on GSM8K and 51.7\% on MATH, outperforming all open-source models within the 7B to 70B parameter range and surpassing most closed-source competitors as well. 2) Notably, DeepSeekMath-RL 7B undergoes instruction tuning exclusively with chain-of-thought-formatted data from GSM8K and MATH, using DeepSeekMath-Instruct 7B as its starting point. Despite the limited diversity of its fine-tuning data, the model demonstrates improved performance over DeepSeekMath-Instruct 7B across all evaluation benchmarks, highlighting the substantial benefits conferred by reinforcement learning.

\section{Discussion}
In this section, we detail the insights and conclusions drawn from our pre-training and RL-based investigations.

\subsection{Lessons Learnt in Pre-Training}

Here, we summarize the key takeaways from our pre-training phase. Unless indicated otherwise, the experiments described here utilize the training configurations established in Section \ref{sec:quality-policy}. When mentioning the DeepSeekMath Corpus in this context, we refer specifically to the 89B-token dataset compiled in the second round of data gathering.

\subsubsection{Code Training Benefits Mathematical Reasoning}

The prevailing, though previously unproven, assumption is that exposure to code during training can enhance a model’s reasoning abilities. Our findings offer partial empirical support for this claim within the realm of mathematics: training on code demonstrably strengthens a model’s mathematical reasoning capacity, both in tasks that require tool usage and those that do not.

To evaluate the influence of code pre-training on mathematical reasoning, we designed experiments involving both a two-stage training regimen and a single-stage training baseline:

\textbf{Two-Stage Training}

\begin{itemize}[topsep=0pt]
    \item \textbf{Code Training for 400B Tokens $\rightarrow$ Math Training for 150B Tokens}: The DeepSeek-LLM 1.3B model undergoes initial pretraining on 400B code tokens, succeeded by additional training with 150B math tokens;
    \item \textbf{General Training for 400B Tokens $\rightarrow$ Math Training for 150B Tokens}: For comparison, we conduct an alternative training strategy where, instead of code tokens, general tokens (sampled from a comprehensive corpus developed by DeepSeek-AI) are used in the first phase. This serves to analyze whether code tokens provide unique benefits for mathematical reasoning, relative to a general token baseline.
\end{itemize}

\textbf{One-Stage Training}

\begin{itemize}[topsep=0pt]
    \item \textbf{Math Training for 150B Tokens}: DeepSeek-LLM 1.3B is directly trained on 150B math tokens only;
    \item \textbf{Training on a mixture of 400B Code Tokens and 150B Math Tokens}: Noting that sequential math training after code pretraining diminishes coding abilities, we further examine if a concurrent mixture of code and math tokens in a single-stage setting sustains or improves both mathematical reasoning and coding skill, and if this integration prevents the loss commonly referred to as catastrophic forgetting.
\end{itemize}

\paragraph{Results}
The performance outcomes for each training configuration are reported in Table \ref{tab:code-math} and Table \ref{tab:code-math-general-eval}.

Incorporating code tokens in training meaningfully aids program-based mathematical reasoning in both the sequential (two-stage) and integrated (one-stage) paradigms. Table \ref{tab:code-math} illustrates that, when applying the two-stage approach, pretraining solely with code tokens already produces substantial advancements in solving GSM8K and MATH questions using Python scripts, while further improvements are observed when math tokens are introduced in the second stage. Notably, in the one-stage scheme, a joint mixture of code and math data lessens the adverse effects of catastrophic forgetting observed in staged approaches, and this blended strategy results in continued gains in both coding capabilities (see Table \ref{tab:code-math-general-eval}) and program-facilitated mathematical problem solving (Table \ref{tab:code-math}).

The presence of code data likewise elevates mathematical reasoning tasks that do not involve external tool usage. Within the two-stage regimen, pretraining with code already provides a moderate positive shift, and it further accelerates the progress achieved through subsequent math-focused training, ultimately delivering the highest overall scores. On the contrary, a simultaneous, one-stage approach using both code and math tokens appears detrimental to tool-free mathematical reasoning—likely because DeepSeek-LLM 1.3B, given its relatively modest model size, may lack the required capacity to comprehensively absorb and unify the diverse information from both code and math sources during a single training phase.

\subsubsection{Limited Effectiveness of ArXiv Papers in Enhancing Mathematical Reasoning}

ArXiv papers are frequently utilized as a source within mathematical pre-training datasets \citep{minerva,gpt-f,llemma,mathpile}. Nevertheless, there has been little systematic assessment of their contribution to mathematical reasoning abilities. Interestingly, our findings suggest that incorporating arXiv papers does not substantially improve mathematical reasoning. Our evaluation covers models of varying sizes—including DeepSeek-LLM 1.3B and DeepSeek-Coder-Base-v1.5 7B \citep{deepseek-coder}—and leverages distinct processing approaches for the arXiv corpus:

\begin{itemize}[topsep=0pt]
    \item \textbf{MathPile} \citep{mathpile}: a dataset containing 8.9B tokens, curated through data cleaning and heuristic filtering, where scientific arXiv papers make up over 85\%;
    \item \textbf{ArXiv-RedPajama} \citep{redpajama}: a dataset comprised of all arXiv LaTeX files, stripped of preambles, comments, macros, and bibliographies, yielding 28.0B tokens.
\end{itemize}

For our experimental procedure, DeepSeek-LLM 1.3B is trained on each arXiv corpus for 150B tokens, and DeepSeek-Coder-Base-v1.5 7B for 40B tokens. Our results indicate that exposure to arXiv data yields negligible or even adverse effects on the models' performance in mathematical reasoning tasks. Across benchmarks of diverse complexity—ranging from quantitative reasoning datasets such as GSM8K and MATH (Table \ref{tab:arxiv-cot}), to multiple-choice exams like MMLU-STEM (Table \ref{tab:arxiv-cot}), and formal mathematics evaluations such as miniF2F (Table \ref{tab:arxiv_atp})—there is little evidence of improvement.

Despite these observations, it is important to acknowledge several caveats. Specifically, our analysis does not address:

\begin{itemize}[topsep=0pt]
    \item The influence of arXiv-sourced content on mathematical tasks outside the scope of our benchmarks, such as transforming formal theorems into informal language;
    \item How arXiv materials might interact with or complement additional pre-training datasets;
    \item Whether larger-scale models might benefit more from arXiv-derived data.
\end{itemize}

Consequently, these outstanding questions warrant future research efforts.

\subsection{Insights of Reinforcement Learning}
\subsubsection{Towards a Unified Paradigm} This section introduces a general framework to compare and analyze different learning algorithms, including SFT, RFT, DPO, PPO, GRPO. We then conduct empirical analysis to investigate how various components within this framework contribute to learning. In general, the gradient of a training algorithm with respect to the parameter $\theta$ can be formulated as follows:

\begin{equation}
    \nabla_{\theta}\mathcal{J}_{\textcolor{red}{\mathcal{A}}}(\theta) = \mathbb{E}[\underbrace{(q,o) \sim \textcolor{red}{\mathcal{D}}}_{Data \ Source}]\left( \frac{1}{|o|} \sum_{t=1}^{|o|}  \underbrace{GC_{{\mathcal{A}}}(q, o, t, \textcolor{red}{\pi_{{rf}}})}_{Gradient \ Coefficient}  \nabla_{\theta}\log \pi_{\theta}(o_t | q, o_{<t})\right).
\label{eq:objective}
\end{equation}

This formulation is composed of three primary elements: 1) \textit{Data Source $\mathcal{D}$}, which specifies the corpus for model training; 2) \textit{Reward Function $\pi_{{rf}}$}, serving as the basis for the reward feedback during training; and 3) \textit{Algorithm $\mathcal{A}$}, responsible for processing both the data and reward to yield the gradient coefficient $GC$ that ultimately governs the strength and direction of updates for each training instance. Within this unified framework, we review and differentiate several commonly used methods:

\begin{itemize}
    \item \textbf{Supervised Fine-tuning (SFT)}: SFT involves continued training of a pretrained model on datasets curated by humans, where the data directly reflect desirable outputs.
    \item \textbf{Rejection Sampling Fine-tuning (RFT)}: In RFT, after initial SFT, the model is further trained on outputs that have been sampled from the SFT model and subsequently filtered to retain only those with correct responses, thus narrowing the data distribution.
    \item \textbf{Direct Preference Optimization (DPO)}: DPO enhances the SFT model further by training with outputs drawn from the SFT model, using a pairwise loss informed by user preferences to directly optimize for preferred generations.
    \item \textbf{Online Rejection Sampling Fine-tuning (Online RFT)}: Unlike standard RFT, Online RFT updates the policy model in real time, sampling outputs and fine-tuning on-the-fly with data drawn from the currently updated model, rather than relying solely on the static SFT model.
    \item \textbf{PPO/GRPO}: In this approach, the policy model is initialized using the SFT model and is iteratively improved using reinforcement learning, drawing new samples from the evolving policy itself and updating parameters based on the corresponding reward signals.
\end{itemize}

We present an overview of these methods' components in Table \ref{tab:data_policy}. For a comprehensive derivation, please consult Appendix \ref{app:analysis-rl}.

\paragraph{Discussion on Data Source} We classify data sources into two types: online sampling and offline sampling. Online sampling implies that the training data originates from the ongoing exploration results of the current policy model, whereas offline sampling indicates reliance on data generated from the initial SFT model’s outputs. Both RFT and DPO utilize the offline approach, whereas Online RFT and GRPO employ online sampling.

As depicted in Figure \ref{fig:rl-analysis}, our observations reveal that Online RFT achieves notably higher performance than RFT across two evaluation benchmarks. In the initial phases of training, Online RFT shows results on par with RFT; however, as training progresses, Online RFT attains a clear and growing lead, emphasizing the benefits of online data collection. This observation is expected, since at the onset, the actor closely resembles the SFT model, leading to minimal differences in sampled data. As training advances, divergences between the actor's outputs and those from the initial model increase, making the advantages of online, policy-driven sampling more pronounced.

\paragraph{Discussion on Gradient Coefficient} The algorithms convert the reward signal into a gradient coefficient for model parameter updates. In our experiments, we categorize the reward function into 'Rule' and 'Model'. The ‘Rule’ method assesses response quality by answer correctness, while the ‘Model’ strategy involves training a reward model—using the rule-based judgments as training data—to assign scores to each response. Equations \ref{eq:GC-RFT} and \ref{eq:GC-GRPO} underline an essential distinction between GRPO and Online RFT: GRPO uniquely modulates its gradient coefficient according to the reward values issued by the reward model, thereby allowing for finer-grained promotion or demotion of responses in line with their reward magnitude. Conversely, Online RFT lacks this specificity: it does not negatively reinforce incorrect responses and applies the same reinforcement intensity to all correctly answered responses.

As illustrated in Figure \ref{fig:rl-analysis}, GRPO outperforms online RFT, underscoring the advantage of dynamically modifying positive and negative gradient coefficients. Moreover, GRPO+PS yields better results than GRPO+OS, suggesting that leveraging detailed, step-sensitive gradient coefficients is beneficial. Additionally, we examine iterative RL by conducting two sequential iterations in our experiments. Figure \ref{fig:iter-rl} demonstrates that iterative RL produces substantial performance gains, particularly evident during the initial iteration.

\subsubsection{Why RL Works?} In this study, reinforcement learning is applied to a subset of instruction tuning data, resulting in considerable improvements over the base instruction-tuned model. To delve into the mechanisms behind RL's effectiveness, we compare the Pass@K and Maj@K accuracies for both Instruct and RL-enhanced models across two benchmark datasets. Figure \ref{fig:maj-pass} shows that RL leads to marked improvements in Maj@K but does not increase Pass@K. These observations imply that RL contributes to overall robustness of the output distribution, meaning \textbf{the benefit primarily lies in elevating the ranking of correct answers within the TopK outputs, rather than advancing core problem-solving ability.} In a similar vein, \citep{wang2023making} describe a \textbf{misalignment problem} in reasoning within SFT models, and note that preference alignment strategies \citep{yuan2023rrhf,song2023preference,wang2023making} are effective for enhancing reasoning capabilities.

\subsubsection{How to Achieve More Effective RL?}
\label{sec:effec_rl}
Our findings show that RL delivers significant benefits in mathematical reasoning scenarios. We also introduce a unified perspective for various prominent training strategies, framing all approaches as either direct applications or simplifications of RL methods. According to Equation \ref{eq:objective}, three essential elements are present: Data Source, Algorithm, and Reward Function. We outline prospective research avenues pertaining to each component.

\paragraph{Data Source} The data source serves as the foundational input for any training paradigm. Specifically in RL contexts, we define the data source as consisting of unlabeled questions accompanied by the outputs generated by the policy model. In the current work, we exclusively utilize questions from the instruction tuning phase and apply a basic nucleus sampling strategy to obtain outputs. This limitation could explain why our RL pipeline predominantly enhances Maj@K. Moving forward, we plan to evaluate our RL framework using out-of-distribution prompts and to incorporate \textbf{more advanced sampling (decoding) strategies}, such as those grounded in tree-search methodologies \citep{yao2023tree}. Additionally, advances in \textbf{efficient inference techniques} \citep{xia-etal-2023-speculative,leviathan2023fast,kwon2023efficient,xia2024unlocking}, which determine how efficiently policy models explore, will continue to be of critical importance.

\paragraph{Algorithms} The algorithmic component is responsible for transforming both the data and the reward signals into gradient coefficients, thereby facilitating the adjustment of the model parameters. As indicated by Equation \ref{eq:objective}, contemporary approaches tend to \textbf{TRUST} the feedback from the reward function almost entirely, leveraging it to modulate the conditional probability of specific tokens. Nonetheless, the infallibility of the reward signal cannot be guaranteed, particularly for highly intricate tasks. For instance, the PRM800K dataset \citep{lightman2023let}, although meticulously curated by expert annotators, is still afflicted by roughly 20\% mislabeling\footnote{\url{https://github.com/openai/prm800k/issues/12\#issuecomment-1728491852}}. To address this shortcoming, we focus on reinforcement learning algorithms capable of operating robustly under noisy or unreliable reward feedback. In our view, the emergence of \textbf{WEAK-TO-STRONG} \citep{burns2023weak} alignment strategies promises to revolutionize existing learning algorithms at a foundational level.

\paragraph{Reward Function} The reward function serves as the principal provider of the training signal in RL settings, commonly instantiated as a neural reward model. We identify three pivotal research avenues for such reward models: 1) \textbf{Enhancing the generalization capacity of reward models.} To effectively improve large language models, reward models must generalize well to previously unseen queries and more sophisticated decoding outputs, lest RL merely reinforces the prevailing distribution rather than advancing intrinsic model competencies; 2) \textbf{Capturing and representing the uncertainty inherent in reward models.} Modeling this uncertainty could be key in connecting underperforming reward models to weak-to-strong algorithmic frameworks; 3) \textbf{Constructing high-fidelity process reward models efficiently} to supply nuanced, step-level training feedback for complex reasoning tasks \citep{lightman2023let,wang2023math}.

\section{Conclusion, Limitation, and Future Work}

We introduce DeepSeekMath, which establishes new benchmarks for open-source models on the competition-level MATH dataset and attains results close to those of proprietary models. DeepSeekMath~builds on DeepSeek-Coder-v1.5 7B as its foundation and is continually trained for 500B tokens; notably, the dataset incorporates 120B math tokens derived from Common Crawl. Through comprehensive ablation experiments, we find that web-based resources serve as a rich source of quality mathematical content, whereas arXiv data does not yield the expected benefits. We propose Group Relative Policy Optimization (GRPO), an adaptation of Proximal Policy Optimization (PPO), which not only enhances mathematical reasoning proficiency but also reduces memory demands. Experimental evidence indicates that GRPO continues to deliver improvements even when DeepSeekMath-Instruct 7B attains strong benchmark scores. Furthermore, we develop a unified framework to interpret a spectrum of methodologies and outline promising avenues for advancing reinforcement learning approaches.

Despite DeepSeekMath's~noteworthy performance on quantitative reasoning tasks, its abilities in geometric and theorem-proving contexts lag behind those of closed-source systems. For example, during initial testing, the model was unable to solve problems involving triangles and ellipses—an outcome that may point to selection biases within the pre-training and fine-tuning datasets. Additionally, model size limitations result in weaker few-shot learning compared to GPT-4; while GPT-4's results improve with additional exemplars, DeepSeekMath~exhibits comparable accuracy in both zero-shot and few-shot settings. In subsequent work, we plan to refine our data selection mechanisms to build an even higher-quality pre-training dataset. Moreover, we aim to investigate directions (see Section \ref{sec:effec_rl}) for enhancing the effectiveness of reinforcement learning in large language models.

\end{document}